\documentclass[a4paper,12pt,twoside]{memoir}

% Castellano
\usepackage[spanish,es-tabla]{babel}
\selectlanguage{spanish}
\usepackage[utf8]{inputenc}
\usepackage[T1]{fontenc}
\usepackage{lmodern} % scalable font
\usepackage{microtype}
\usepackage{placeins}

\RequirePackage{booktabs}
\RequirePackage[table]{xcolor}
\RequirePackage{xtab}
\RequirePackage{multirow}

% Links
\PassOptionsToPackage{hyphens}{url}\usepackage[colorlinks]{hyperref}
\hypersetup{
	allcolors = {red}
}

% Ecuaciones
\usepackage{amsmath}

% Rutas de fichero / paquete
\newcommand{\ruta}[1]{{\sffamily #1}}

% Párrafos
\nonzeroparskip

% Huérfanas y viudas
\widowpenalty100000
\clubpenalty100000

% Evitar solapes en el header
\nouppercaseheads

% Imagenes
\usepackage{graphicx}
\newcommand{\imagen}[2]{
	\begin{figure}[!h]
		\centering
		\includegraphics[width=0.9\textwidth]{#1}
		\caption{#2}\label{fig:#1}
	\end{figure}
	\FloatBarrier
}

\newcommand{\imagenflotante}[2]{
	\begin{figure}%[!h]
		\centering
		\includegraphics[width=0.9\textwidth]{#1}
		\caption{#2}\label{fig:#1}
	\end{figure}
}



% El comando \figura nos permite insertar figuras comodamente, y utilizando
% siempre el mismo formato. Los parametros son:
% 1 -> Porcentaje del ancho de página que ocupará la figura (de 0 a 1)
% 2 --> Fichero de la imagen
% 3 --> Texto a pie de imagen
% 4 --> Etiqueta (label) para referencias
% 5 --> Opciones que queramos pasarle al \includegraphics
% 6 --> Opciones de posicionamiento a pasarle a \begin{figure}
\newcommand{\figuraConPosicion}[6]{%
  \setlength{\anchoFloat}{#1\textwidth}%
  \addtolength{\anchoFloat}{-4\fboxsep}%
  \setlength{\anchoFigura}{\anchoFloat}%
  \begin{figure}[#6]
    \begin{center}%
      \Ovalbox{%
        \begin{minipage}{\anchoFloat}%
          \begin{center}%
            \includegraphics[width=\anchoFigura,#5]{#2}%
            \caption{#3}%
            \label{#4}%
          \end{center}%
        \end{minipage}
      }%
    \end{center}%
  \end{figure}%
}

%
% Comando para incluir imágenes en formato apaisado (sin marco).
\newcommand{\figuraApaisadaSinMarco}[5]{%
  \begin{figure}%
    \begin{center}%
    \includegraphics[angle=90,height=#1\textheight,#5]{#2}%
    \caption{#3}%
    \label{#4}%
    \end{center}%
  \end{figure}%
}
% Para las tablas
\newcommand{\otoprule}{\midrule [\heavyrulewidth]}
%
% Nuevo comando para tablas pequeñas (menos de una página).
\newcommand{\tablaSmall}[5]{%
 \begin{table}
  \begin{center}
   \rowcolors {2}{gray!35}{}
   \begin{tabular}{#2}
    \toprule
    #4
    \otoprule
    #5
    \bottomrule
   \end{tabular}
   \caption{#1}
   \label{tabla:#3}
  \end{center}
 \end{table}
}

%
%Para el float H de tablaSmallSinColores
\usepackage{float}

%
% Nuevo comando para tablas pequeñas (menos de una página).
\newcommand{\tablaSmallSinColores}[5]{%
 \begin{table}[H]
  \begin{center}
   \begin{tabular}{#2}
    \toprule
    #4
    \otoprule
    #5
    \bottomrule
   \end{tabular}
   \caption{#1}
   \label{tabla:#3}
  \end{center}
 \end{table}
}

\newcommand{\tablaApaisadaSmall}[5]{%
\begin{landscape}
  \begin{table}
   \begin{center}
    \rowcolors {2}{gray!35}{}
    \begin{tabular}{#2}
     \toprule
     #4
     \otoprule
     #5
     \bottomrule
    \end{tabular}
    \caption{#1}
    \label{tabla:#3}
   \end{center}
  \end{table}
\end{landscape}
}

%
% Nuevo comando para tablas grandes con cabecera y filas alternas coloreadas en gris.
\newcommand{\tabla}[6]{%
  \begin{center}
    \tablefirsthead{
      \toprule
      #5
      \otoprule
    }
    \tablehead{
      \multicolumn{#3}{l}{\small\sl continúa desde la página anterior}\\
      \toprule
      #5
      \otoprule
    }
    \tabletail{
      \hline
      \multicolumn{#3}{r}{\small\sl continúa en la página siguiente}\\
    }
    \tablelasttail{
      \hline
    }
    \bottomcaption{#1}
    \rowcolors {2}{gray!35}{}
    \begin{xtabular}{#2}
      #6
      \bottomrule
    \end{xtabular}
    \label{tabla:#4}
  \end{center}
}

%
% Nuevo comando para tablas grandes con cabecera.
\newcommand{\tablaSinColores}[6]{%
  \begin{center}
    \tablefirsthead{
      \toprule
      #5
      \otoprule
    }
    \tablehead{
      \multicolumn{#3}{l}{\small\sl continúa desde la página anterior}\\
      \toprule
      #5
      \otoprule
    }
    \tabletail{
      \hline
      \multicolumn{#3}{r}{\small\sl continúa en la página siguiente}\\
    }
    \tablelasttail{
      \hline
    }
    \bottomcaption{#1}
    \begin{xtabular}{#2}
      #6
      \bottomrule
    \end{xtabular}
    \label{tabla:#4}
  \end{center}
}

%
% Nuevo comando para tablas grandes sin cabecera.
\newcommand{\tablaSinCabecera}[5]{%
  \begin{center}
    \tablefirsthead{
      \toprule
    }
    \tablehead{
      \multicolumn{#3}{l}{\small\sl continúa desde la página anterior}\\
      \hline
    }
    \tabletail{
      \hline
      \multicolumn{#3}{r}{\small\sl continúa en la página siguiente}\\
    }
    \tablelasttail{
      \hline
    }
    \bottomcaption{#1}
  \begin{xtabular}{#2}
    #5
   \bottomrule
  \end{xtabular}
  \label{tabla:#4}
  \end{center}
}



\definecolor{cgoLight}{HTML}{EEEEEE}
\definecolor{cgoExtralight}{HTML}{FFFFFF}

%
% Nuevo comando para tablas grandes sin cabecera.
\newcommand{\tablaSinCabeceraConBandas}[5]{%
  \begin{center}
    \tablefirsthead{
      \toprule
    }
    \tablehead{
      \multicolumn{#3}{l}{\small\sl continúa desde la página anterior}\\
      \hline
    }
    \tabletail{
      \hline
      \multicolumn{#3}{r}{\small\sl continúa en la página siguiente}\\
    }
    \tablelasttail{
      \hline
    }
    \bottomcaption{#1}
    \rowcolors[]{1}{cgoExtralight}{cgoLight}

  \begin{xtabular}{#2}
    #5
   \bottomrule
  \end{xtabular}
  \label{tabla:#4}
  \end{center}
}




\graphicspath{ {./img/} }

% Capítulos
\chapterstyle{bianchi}
\newcommand{\capitulo}[2]{
	\setcounter{chapter}{#1}
	\setcounter{section}{0}
	\setcounter{figure}{0}
	\setcounter{table}{0}
	\chapter*{#2}
	\addcontentsline{toc}{chapter}{#2}
	\markboth{#2}{#2}
}

% Apéndices
\renewcommand{\appendixname}{Apéndice}
\renewcommand*\cftappendixname{\appendixname}

\newcommand{\apendice}[1]{
	%\renewcommand{\thechapter}{A}
	\chapter{#1}
}

\renewcommand*\cftappendixname{\appendixname\ }

% Formato de portada
\makeatletter
\usepackage{xcolor}
\newcommand{\tutor}[1]{\def\@tutor{#1}}
\newcommand{\course}[1]{\def\@course{#1}}
\definecolor{cpardoBox}{HTML}{E6E6FF}
\def\maketitle{
  \null
  \thispagestyle{empty}
  % Cabecera ----------------
\noindent\includegraphics[width=\textwidth]{cabecera}\vspace{1cm}%
  \vfill
  % Título proyecto y escudo informática ----------------
  \colorbox{cpardoBox}{%
    \begin{minipage}{.8\textwidth}
      \vspace{.5cm}\Large
      \begin{center}
      \textbf{TFG del Grado en Ingeniería Informática}\vspace{.6cm}\\
      \textbf{\LARGE\@title{}}
      \end{center}
      \vspace{.2cm}
    \end{minipage}

  }%
  \hfill\begin{minipage}{.20\textwidth}
    \includegraphics[width=\textwidth]{escudoInfor}
  \end{minipage}
  \vfill
  % Datos de alumno, curso y tutores ------------------
  \begin{center}%
  {%
    \noindent\LARGE
    Presentado por \@author{}\\ 
    en Universidad de Burgos --- \@date{}\\
    Tutor: \@tutor{}\\
  }%
  \end{center}%
  \null
  \cleardoublepage
  }
\makeatother


% Datos de portada
\title{título del TFG \\Documentación Técnica}
\author{nombre alumno}
\tutor{nombre tutor}
\date{\today}

\begin{document}

\maketitle



\cleardoublepage



%%%%%%%%%%%%%%%%%%%%%%%%%%%%%%%%%%%%%%%%%%%%%%%%%%%%%%%%%%%%%%%%%%%%%%%%%%%%%%%%%%%%%%%%



\frontmatter


\clearpage

% Indices
\tableofcontents

\clearpage

\listoffigures

\clearpage

\listoftables

\clearpage

\mainmatter

\appendix

\apendice{Plan de Proyecto Software}

\section{Introducción}
Este primer apartado de los anexos se ha dedicado a la explicación del
desarrollo del proyecto, recogido en planificación temporal, y a la viabilidad
del proyecto, dividida en viabilidad económica donde se hará una aproximación de los costes que conlleva el proyecto y viabilidad legal donde se debe
tener en cuenta la legislación que puede afectar al proyecto.

\section{Planificación temporal}
Antes de comenzar el proyecto se decidió utilizar una metodología ágil para
gestionarlo. Se han planificado sprints de dos semanas de duración.
Al final de cada sprint se ha hecho revisión para comprobar si se habían
conseguido los objetivos marcados al inicio y plantear las tareas a
realizar para el nuevo sprint.
Para la gestión del proyecto se han utilizado GitHub y Zube. El
tablero que provee Zube ha permitido que visualmente la organización
del proyecto haya sido más sencilla. Desde github se han vinculado los ``commits'' 
con cada tarea para hacer más fácil su seguimiento.

\subsection{Sprint 0 - del 2 de Octubre de 2023 al 15 de Octubre de 2023}
Este sprint 0 se ha dedicado a tareas de preparación.
Se ha creado el repositorio en Github, un proyecto en Zube y se ha instalado 
un entorno Latex para la documentación.
Por otro lado me he dedicado a leer la documentación relacionada y 
realizar los cursos recomendados.
\imagen{tareas_sprint_0}{Tareas sprint 0}
\imagen{burndown_sprint_0}{Burndown chart sprint 0}
\subsection{Sprint 1 - del 16 de Octubre de 2023 al 29 de Octubre de 2023}
En el sprint 1 he empezado a definir el catálogo de especificaciones y empezar a 
pensar cómo implementar tanto el prototipo como la aplicación de usuario.
También he buscado información sobre trabajos relacionados.
\imagen{tareas_sprint_1}{Tareas sprint 1}
\imagen{burndown_sprint_1}{Burndown chart sprint 1}
\subsection{Sprint 2 - del 30 de Octubre de 2023 al 12 de Noviembre de 2023}
En este segundo sprint por un lado he empezado a definir la arquitectura de la aplicación 
y por otro he empezado con el prototipo en collab. Para ello primero he tenido que buscar 
un dataset que nos valga y un modelo que se pueda usar en collab.
Después de mucho buscar decidí usar un modelo llama (con llama.cpp) y un dataset de hugging face 
para hacerlo más portable.
Si la máquina asignada tiene tarjeta gráfica, se usará.

\imagen{tareas_sprint_2}{Tareas sprint 2}
\imagen{burndown_sprint_2}{Burndown chart sprint 2}
\subsection{Sprint 3 - del 13 de Noviembre de 2023 al 26 de Noviembre de 2023}
En este sprint empecé a definir los casos de uso, el interfaz de usuario y empezar a pensar 
en el modelo de base de datos. En este momento decidí usar un ORM para hacer más flexible la 
aplicación e independiente del SGBD. Así sería más fácil cambiar de SGBD en caso de encontrar 
algún problema. En este momento también decidí empezar a usar postgres y docker por simplicidad.

\imagen{tareas_sprint_3}{Tareas sprint 3}
\imagen{burndown_sprint_3}{Burndown chart sprint 3}
\subsection{Sprint 4 - del 27 de Noviembre de 2023 al 10 de Diciembre de 2023}
En este sprint ya empecé más en serio con el desarrollo de la aplicación de usuario con
la implementación de la pantalla de login, la pantalla de administración y la pantalla 
de ejecución del proceso.
Por otro lado también empecé a con la implentación de la base de datos, el cliente OpenAI y 
la elección de un hosting.
\imagen{tareas_sprint_4}{Tareas sprint 4}
\imagen{burndown_sprint_4}{Burndown chart sprint 4}
\subsection{Sprint 5 - del 11 de Diciembre de 2023 al 24 de Diciembre de 2023}
En este sprint se empieza a probar la aplicación, se 
resuelven los errores encontrados y 
se empieza con la documentación técnica.
\imagen{tareas_sprint_5}{Tareas sprint 5}
\imagen{burndown_sprint_5}{Burndown chart sprint 5}
\subsection{Sprint 6 - del 25 de Diciembre de 2023 al 7 de Enero de 2024}
A raíz de una reunión con el tutor se propone un cambio en el enfoque de la aplicación y 
se crea una nueva pantalla de análisis de reseñas.
También se empieza con el estudio de viabilidad.
\imagen{tareas_sprint_6}{Tareas sprint 6}
\imagen{burndown_sprint_6}{Burndown chart sprint 6}
\subsection{Sprint 7 - del 7 de Enero de 2024 al 21 de Enero de 2024}
Se continua con las pruebas, y se corrige la memoria con las sugerencias de los tutores.
Se avanza con la memoria y se implementa la traducción de la aplicación.
Se configura Codacy para las pruebas de calidad del código.
\imagen{tareas_sprint_7}{Tareas sprint 7}
\imagen{burndown_sprint_7}{Burndown chart sprint 7}

\section{Estudio de viabilidad}

\subsection{Viabilidad Técnica}

Las tecnologías utilizadas en este trabajo se emplean ampliamente en muchos proyectos de software, y 
tienen un amplio soporte por parte de su comunidad de desarrollo. Son conocidas por su eficiencia y 
flexibilidad en el desarrollo web.\\
Basándonos en los requisitos del proyecto las tecnologías elegidas para el desarrollo del trabajo 
resultan adecuadas para su correcto funcionamiento.\\
La mayor complejidad del trabajo se encuentra en el modelo LLM que se gestiona desde una instancia ajena al 
resto del proyecto de desarrollo lo que nos facilita su gestión.\\
Para el caso concreto de OpenAI resulta muy útil su uso pero hay que tener en cuenta los costes asociados 
así como sus políticas de uso.

\subsection{Viabilidad económica}
La viabilidad económica de un proyecto de software construido con las tecnologías 
mencionadas dependerá de varios factores, incluyendo el alcance y los objetivos del proyecto, 
el mercado al que se dirige y la estrategia de monetización.

\subsubsection{Desarrollo}
En este apartado vamos a calcular el coste de desarrollo del proyecto.
Por un lado tenemos el coste del desarrollador.\\
Se han empleado alrededor de 4 meses para completar el trabajo, vamos a considerar 
un 50% por ciento de dedicación.\\
Por lo tanto tenemos unas 20 horas semanales durante 16 semanas.
El sueldo medio por hora de un desarrollador en España es de unos \href{https://es.talent.com}{15 euros} 
Así el salario total ha sido: 20*16*15 = 4800 euros
A esto hay que agregar el coste de la cotización a la seguridad social.
En estos momentos la cotización por parte de la empresa se distribuye en 
\footnote{https://www.seg-social.es/wps/portal/wss/internet/Trabajadores/CotizacionRecaudacionTrabajadores/36537}:
\begin{itemize}
    \item Contigencias: 23,6%
    \item Desempleo: 5,5%
    \item FOGASA: 0,20
    \item Formación profesional (FP): 0,60%
    \item Mecanismo de equidad intergeneracional (MEI): 0,5%
\end{itemize}
Asi el coste total será:

4800(salario) + 1132,8 (contigencias) + 264 (desempleo) + 9,6 (FOGASA) + 28,8 (FP) + 24 (MEI) = 6259,2

De la misma forma se valora el coste de asesoría (tutores), contando 2 horas a la
 semana con un coste de 40 euros la hora. 
En este caso el salario total sería:
40*2*16*2 = 2560

Y el coste total: 

2560(salario) + 604,16(contigencias) +  140 (desempleo) + 5,12 (FOGASA) + 15,36 (FP) + 12,8 (MEI) = 3337,44

\subsubsection{Software}
Todo el software utilizado en este proyecto es gratuito para aplicaciones no comerciales pero tienen un 
coste para aplicaciones comerciales excepto OpenAI que siempre tiene un coste por uso.

\begin{tabular}{r r}
    \textbf{Software} & \textbf{Coste}\\
    Python & Gratuito \\
    node.js & Gratuito\\
    Docker & Gratuito\\
    Vue & Gratuito\\
    Postgres & Gratuito\\
    OpenAI & Por uso\\
\end{tabular}


Para calcular el coste de OpenAI tenemos que tener en cuenta que éste se calcula según el 
número de \href{https://platform.openai.com/tokenizer}{tokens} enviados/recibidos.\footnote{fuente: https://openai.com/pricing}

\begin{tabular}{r r r}
    \textbf{Model}&\textbf{Input}&\textbf{Output}\\
    gpt-3.5-turbo-1106 & \$0.0010 / 1K tokens & \$0.0020 / 1K tokens\\
    gpt-3.5-turbo-instruct & \$0.0015 / 1K tokens & \$0.0020 / 1K tokens\\        
\end{tabular}

1000 tokens son alrededor de 750 palabras.

\subsubsection{Hardware}
Acerca del hardware utilizado en el proyecto, se ha utilizado un alojamiento llamado \href{https://contabo.com/en/}{Contabo}.\\
Este alojamiento tiene un coste de 4,5 euros al mes (+IVA = 5,45€). 
Ese precio incluye un máquina virtual linux (Ubuntu),
con 8Gb de Memoria RAM y 200Gb de disco duro.\\
Esta configuración es suficiente para el proyecto realizado.

El resto de hardware utilizado se considera ya amortizado en proyectos 
pasados y por tanto no se va a contar para este.


\subsubsection{Otros}

Otros costes asociados a este desarrollo serían:
Conexión a internet (40€ al mes)
El alquiler del dominio, que ha costado 1€ al año en \href{https://www.arsys.es/}{arsys}.
El certificado SSL se ha obtenido gratuitamente de \href{https://letsencrypt.org/}{letsencrypt}.

\subsubsection{Total}

\begin{tabular}{l r}
    \textbf{Concepto} & \textbf{Coste}\\
    Desarollo &  9596,54\\
    Software & 10€\\
    Hardware & 20€\\
    Otros & 161€\\
    Total & 9787,54\\
\end{tabular}

\subsection{Viabilidad legal}

Para estudiar la viabilidad legal del proyecto vamos a analizar las licencias de las
herramientas utilizadas.

\begin{tabular}{l l}
    \textbf{Software} & \textbf{Licencia}\\
    Python & PSF (compatible con GPL) \\
    Flask & BSD\\
    SQLAlchemy & MIT\\
    Github & GNU\\
    node.js & MIT\\
    Docker & Apache 2.0\\
    Vue & MIT\\
    Postgres & Licencia Postgres\\
    OpenAI & Propietaria\\
    Visual Studio Code & MIT \\
\end{tabular}

Como vemos todas las herramientas utilizadas durante el desarrollo son Open Source, 
excepto OpenAI.
\apendice{Especificación de Requisitos}
\section{Introducción}
\section{Objetivos generales}

Crear una aplicación que permita realizar análisis de sentimiento a textos extraídos de reseñas de Amazon.
Almacenar dicha información en una base de datos para que sea fácilmente accesible.
Mostrar dicha información con gráficos y tablas para poder analizarlos.

Crear un manual de usuario. 

\section{Catálogo de requisitos}
\subsection{Requisitos funcionales}
\begin{itemize}
	\tightlist
	\item \textbf{RF1 - Registro de usuarios}: La aplicación debe permitir dar de alta usuarios y asignarles un rol.
	\item \textbf{RF2 - Login de usuarios}: La aplicación debe ser capaz de dadas las credenciales de un usuario permitirle acceder a las funcionalidades. 
	\item \textbf{RF3 - Información de la aplicación}: El usuario debe poder obtener información sobre el funcionamiento de la aplicación. 
	\item \textbf{RF4 - Analizar textos}: La aplicación debe ser capaz de analizar textos.
	\begin{itemize}
		\item \textbf{RF4.1 - Cargar reseñas de un archivo}: La aplicación debe poder cargar las reseñas de un archivo de texto
		\item \textbf{RF4.2 - Puntuar los textos}: La aplicación debe realizar el análisis del texto obtenido y devolver un valor entero entre 1 y 5 dependiendo del sentimiento del mismo, siendo 1 un sentimiento negativo y 5 un sentimiento positivo.     
		\item \textbf{RF4.3 - Almacenar resultados}: La aplicación debe almacenar las puntuaciones junto con el texto analizado, fecha en la que se escribió y un identificador. 
	\end{itemize}
	\item \textbf{RF5 - Buscar resultados por identificador}: El usuario debe poder buscar los resultados de un identificador con unas fechas seleccionadas. 
	\begin{itemize}
		\item \textbf{RF5.1 - Buscar por producto}: El usuario debe poder introducir el nombre de un producto y mostrar las reseñas del mismo. 
		\item \textbf{RF5.2 - Buscar por usuario}: El usuario debe poder buscar un usuario y mostrar la información de las reseñas que ha creado. 
		\item \textbf{RF5.3 - Buscar por texto}: El usuario debe poder buscar una palabra para ver los resultados del análisis de los textos en los que está esa palabra.
	\end{itemize}
	\item \textbf{RF6 - Matriz de confusión}: Debe mostrar la matriz de confusión para los datos seleccionados, comparando los resultados obtenidos y la calificación dada por el usuario.
	\item \textbf{RF7 - Configurar idioma}: El usuario debe poder cambiar el idioma de la aplicación. 
\end{itemize}      
\subsection{Requisitos no funcionales}
\begin{itemize}
	\item \textbf{RNF1 - Usabilidad}: La aplicación debe ser amigable para que la experiencia del usuario sea lo más positiva posible. Además debe poder adaptarse a diferentes formatos de pantalla. 
	\item \textbf{RNF2 - Eficiencia}: El almacenamiento de resultados en la base de datos y la respuesta de la aplicación al navegar entre ventanas, sobre todo, al mostrar los gráficos debe ser lo más rápido posible. 
	\item \textbf{RNF3 - Disponibilidad}: La aplicación debe estar disponible durante el mayor tiempo posible y en la mayoría de lugares. 
	\item \textbf{RNF4 - Escalabilidad}: La aplicación debe estar preparada para soportar nuevos desarrollos que generen mayor cantidad de trabajo. 
	\item \textbf{RNF5 - Confiabilidad}: La aplicación cumplirá la función para la que se ha creado. 
	\item \textbf{RNF6 - Mantenibilidad}: La aplicación debe permitir cambios y correcciones de forma eficaz.
	\item \textbf{RNF7 - Internacionalización}: La aplicación debe soportar varios idiomas.
	\item \textbf{RNF8 - Seguridad}: La aplicación debe soportar usuarios con diferentes roles que tendrán acceso a diferentes funcionalidades.
	
\end{itemize}


\section{Casos de uso}

% Caso de Uso 1 -> Registro de usuarios.
\begin{table}[p]
	\centering
	\begin{tabularx}{\linewidth}{ p{0.21\columnwidth} p{0.71\columnwidth} }
		\toprule
		\textbf{CU-1}    & \textbf{Registro de usuarios}\\
		\toprule
		\textbf{Versión}              & 1.0    \\
		\textbf{Autor}                & Teodoro Ricardo García Sánchez \\
		\textbf{Requisitos asociados} & RF-1 \\
		\textbf{Descripción}          & Un usuario con el rol de administrador podrá crear cuentas de usuario y asignarles un rol. \\
		\textbf{Precondición}         & El usuario tiene que tener el rol de administrador \\
		\textbf{Acciones}             &
		\begin{enumerate}
			\def\labelenumi{\arabic{enumi}.}
			\tightlist
			\item El usuario administrador accede a la interfaz de la aplicación.
			\item Accede a la opción Crear usuarios
			\item Introduce los datos del usuarios
			\item Confirma los cambios.
		\end{enumerate}\\
		\textbf{Postcondición}        & La cuenta de usuario no debe existir con antelación \\
		\textbf{Excepciones}          & 
		\begin{enumerate}
			\item El usuario ya existe
			\item Complejidad de la contraseña no válida
			\item Datos incorrectos
		\end{enumerate}\\
		\textbf{Importancia}          & Alta \\
		\bottomrule
	\end{tabularx}
	\caption{CU-1 Registro de usuarios.}
\end{table}
% Caso de Uso 2 -> Login de usuarios.
\begin{table}[p]
	\centering
	\begin{tabularx}{\linewidth}{ p{0.21\columnwidth} p{0.71\columnwidth} }
		\toprule
		\textbf{CU-2}    & \textbf{Login de usuarios}\\
		\toprule
		\textbf{Versión}              & 1.0    \\
		\textbf{Autor}                & Teodoro Ricardo García Sánchez \\
		\textbf{Requisitos asociados} & RF-2 \\
		\textbf{Descripción}          & El usuario accede a la aplicación tendiendo \\
										acceso a la parte que tengo acceso según su rol. \\
		\textbf{Precondición}         & Usuario registrado \\
		\textbf{Acciones}             &
		\begin{enumerate}
			\def\labelenumi{\arabic{enumi}.}
			\tightlist
			\item Accede a la pantalla de login 
			\item Introduce nombre y contraseña
			\item Pulsa en Aceptar
		\end{enumerate}\\
		\textbf{Postcondición}        & La contraseña debe coincidir \\
		\textbf{Excepciones}          & 
		\begin{enumerate}
			\def\labelenumi{\arabic{enumi}.}
			\tightlist
			\item El usuario no existe 
			\item La contraseña no es válida 
		\end{enumerate}\\
		\textbf{Importancia}          & Alta\\
		\bottomrule
	\end{tabularx}
	\caption{CU-2 Login de usuarios.}
\end{table}
% Caso de Uso 3 -> Información de la aplicación.
\begin{table}[p]
	\centering
	\begin{tabularx}{\linewidth}{ p{0.21\columnwidth} p{0.71\columnwidth} }
		\toprule
		\textbf{CU-3}    & \textbf{Información de la aplicación}\\
		\toprule
		\textbf{Versión}              & 1.0    \\
		\textbf{Autor}                & Teodoro Ricardo García Sánchez \\
		\textbf{Requisitos asociados} & RF-3 \\
		\textbf{Descripción}          & Permite al usuario obtener información sobre la configuración de la aplicación \\
		\textbf{Precondición}         & Usuario logado \\
		\textbf{Acciones}             &
		\begin{enumerate}
			\def\labelenumi{\arabic{enumi}.}
			\tightlist
			\item El usuario accede a la opción información.
		\end{enumerate}\\
		\textbf{Postcondición}        &  \\
		\textbf{Excepciones}          &  \\
		\textbf{Importancia}          & Baja \\
		\bottomrule
	\end{tabularx}
	\caption{CU-3 Información de la aplicación.}
\end{table}
% Caso de Uso 4 -> Cargar reseñas de un archivo.
\begin{table}[p]
	\centering
	\begin{tabularx}{\linewidth}{ p{0.21\columnwidth} p{0.71\columnwidth} }
		\toprule
		\textbf{CU-4}    & \textbf{Cargar reseñas de un archivo}\\
		\toprule
		\textbf{Versión}              & 1.0    \\
		\textbf{Autor}                & Teodoro Ricardo García Sánchez \\
		\textbf{Requisitos asociados} & RF4.1 \\
		\textbf{Descripción}          & El usuario puede cargar reseñas de un archivo \\
		\textbf{Precondición}         & Usuario logado \\
		\textbf{Acciones}             &
		\begin{enumerate}
			\def\labelenumi{\arabic{enumi}.}
			\tightlist
			\item Pasos del CU
			\item Pasos del CU (añadir tantos como sean necesarios)
		\end{enumerate}\\
		\textbf{Postcondición}        & Postcondiciones (podría haber más de una) \\
		\textbf{Excepciones}          & Excepciones \\
		\textbf{Importancia}          & Alta o Media o Baja... \\
		\bottomrule
	\end{tabularx}
	\caption{CU-04 Cargar reseñas de un archivo.}
\end{table}
% Caso de Uso 5 -> Puntuar los textos.
\begin{table}[p]
	\centering
	\begin{tabularx}{\linewidth}{ p{0.21\columnwidth} p{0.71\columnwidth} }
		\toprule
		\textbf{CU-5}    & \textbf{Puntuar los textos}\\
		\toprule
		\textbf{Versión}              & 1.0    \\
		\textbf{Autor}                & Teodoro Ricardo García Sánchez \\
		\textbf{Requisitos asociados} & RF4.2 \\
		\textbf{Descripción}          & La aplicación debe realizar el análisis del texto obtenido
										 y devolver un valor entero entre 1 y 5 dependiendo del sentimiento del mismo, siendo 1 un sentimiento negativo y 5 un sentimiento positivo. \\
		\textbf{Precondición}         &  \\
		\textbf{Acciones}             &
		\begin{enumerate}
			\def\labelenumi{\arabic{enumi}.}
			\tightlist
			\item Pasos del CU
			\item Pasos del CU (añadir tantos como sean necesarios)
		\end{enumerate}\\
		\textbf{Postcondición}        & Postcondiciones (podría haber más de una) \\
		\textbf{Excepciones}          & Excepciones \\
		\textbf{Importancia}          & Alta o Media o Baja... \\
		\bottomrule
	\end{tabularx}
	\caption{CU-5 Puntuar los textos.}
\end{table}
% Caso de Uso 6 -> Almacenar resultados.
\begin{table}[p]
	\centering
	\begin{tabularx}{\linewidth}{ p{0.21\columnwidth} p{0.71\columnwidth} }
		\toprule
		\textbf{CU-6}    & \textbf{Alamacenar resultados}\\
		\toprule
		\textbf{Versión}              & 1.0    \\
		\textbf{Autor}                & Teodoro Ricardo García Sánchez \\
		\textbf{Requisitos asociados} & RF4.3 \\
		\textbf{Descripción}          & La aplicación debe almacenar las puntuaciones junto con el texto analizado, fecha en la que se escribió y un identificador. \\
		\textbf{Precondición}         & Precondiciones (podría haber más de una) \\
		\textbf{Acciones}             &
		\begin{enumerate}
			\def\labelenumi{\arabic{enumi}.}
			\tightlist
			\item Pasos del CU
			\item Pasos del CU (añadir tantos como sean necesarios)
		\end{enumerate}\\
		\textbf{Postcondición}        & Postcondiciones (podría haber más de una) \\
		\textbf{Excepciones}          & Excepciones \\
		\textbf{Importancia}          & Alta o Media o Baja... \\
		\bottomrule
	\end{tabularx}
	\caption{CU-6 Almacenar resultados.}
\end{table}
% Caso de Uso 7 -> Buscar resultados por producto.
\begin{table}[p]
	\centering
	\begin{tabularx}{\linewidth}{ p{0.21\columnwidth} p{0.71\columnwidth} }
		\toprule
		\textbf{CU-7}    & \textbf{Buscar resultado por producto}\\
		\toprule
		\textbf{Versión}              & 1.0    \\
		\textbf{Autor}                & Teodoro Ricardo García Sánchez \\
		\textbf{Requisitos asociados} & RF5.1 \\
		\textbf{Descripción}          & El usuario debe poder introducir el nombre de un producto y mostrar las reseñas del mismo \\
		\textbf{Precondición}         & Precondiciones (podría haber más de una) \\
		\textbf{Acciones}             &
		\begin{enumerate}
			\def\labelenumi{\arabic{enumi}.}
			\tightlist
			\item Pasos del CU
			\item Pasos del CU (añadir tantos como sean necesarios)
		\end{enumerate}\\
		\textbf{Postcondición}        & Postcondiciones (podría haber más de una) \\
		\textbf{Excepciones}          & Excepciones \\
		\textbf{Importancia}          & Alta o Media o Baja... \\
		\bottomrule
	\end{tabularx}
	\caption{CU-7 Buscar resultados por producto.}
\end{table}
% Caso de Uso 8 -> Buscar por nombre de usuario.
\begin{table}[p]
	\centering
	\begin{tabularx}{\linewidth}{ p{0.21\columnwidth} p{0.71\columnwidth} }
		\toprule
		\textbf{CU-8}    & \textbf{Buscar por nombre de usuario}\\
		\toprule
		\textbf{Versión}              & 1.0    \\
		\textbf{Autor}                & Teodoro Ricardo García Sánchez \\
		\textbf{Requisitos asociados} & RF5.2 \\
		\textbf{Descripción}          & Se debe poder buscar un usuario y mostrar la información de las reseñas que ha creado. \\
		\textbf{Precondición}         & Precondiciones (podría haber más de una) \\
		\textbf{Acciones}             &
		\begin{enumerate}
			\def\labelenumi{\arabic{enumi}.}
			\tightlist
			\item Pasos del CU
			\item Pasos del CU (añadir tantos como sean necesarios)
		\end{enumerate}\\
		\textbf{Postcondición}        & Postcondiciones (podría haber más de una) \\
		\textbf{Excepciones}          & Excepciones \\
		\textbf{Importancia}          & Alta o Media o Baja... \\
		\bottomrule
	\end{tabularx}
	\caption{CU-8 Buscar nombre de usuario.}
\end{table}
% Caso de Uso 9 -> Buscar por texto.
\begin{table}[p]
	\centering
	\begin{tabularx}{\linewidth}{ p{0.21\columnwidth} p{0.71\columnwidth} }
		\toprule
		\textbf{CU-9}    & \textbf{Buscar por texto}\\
		\toprule
		\textbf{Versión}              & 1.0    \\
		\textbf{Autor}                & Teodoro Ricardo García Sánchez \\
		\textbf{Requisitos asociados} & RF5.3 \\
		\textbf{Descripción}          & El usuario debe poder buscar una palabra para ver los resultados del análisis de los textos en los que está esa palabra. \\
		\textbf{Precondición}         & Precondiciones (podría haber más de una) \\
		\textbf{Acciones}             &
		\begin{enumerate}
			\def\labelenumi{\arabic{enumi}.}
			\tightlist
			\item Pasos del CU
			\item Pasos del CU (añadir tantos como sean necesarios)
		\end{enumerate}\\
		\textbf{Postcondición}        & Postcondiciones (podría haber más de una) \\
		\textbf{Excepciones}          & Excepciones \\
		\textbf{Importancia}          & Alta o Media o Baja... \\
		\bottomrule
	\end{tabularx}
	\caption{CU-9 Buscar por texto.}
\end{table}
% Caso de Uso 10 -> Matriz de confusión.
\begin{table}[p]
	\centering
	\begin{tabularx}{\linewidth}{ p{0.21\columnwidth} p{0.71\columnwidth} }
		\toprule
		\textbf{CU-8}    & \textbf{Matriz de confusión}\\
		\toprule
		\textbf{Versión}              & 1.0    \\
		\textbf{Autor}                & Teodoro Ricardo García Sánchez \\
		\textbf{Requisitos asociados} & RF-xx, RF-xx \\
		\textbf{Descripción}          & La descripción del CU \\
		\textbf{Precondición}         & Precondiciones (podría haber más de una) \\
		\textbf{Acciones}             &
		\begin{enumerate}
			\def\labelenumi{\arabic{enumi}.}
			\tightlist
			\item Pasos del CU
			\item Pasos del CU (añadir tantos como sean necesarios)
		\end{enumerate}\\
		\textbf{Postcondición}        & Postcondiciones (podría haber más de una) \\
		\textbf{Excepciones}          & Excepciones \\
		\textbf{Importancia}          & Alta o Media o Baja... \\
		\bottomrule
	\end{tabularx}
	\caption{CU-1 Nombre del caso de uso.}
\end{table}

% Caso de Uso 11 -> Consultar Experimentos.
\begin{table}[p]
	\centering
	\begin{tabularx}{\linewidth}{ p{0.21\columnwidth} p{0.71\columnwidth} }
		\toprule
		\textbf{CU-1}    & \textbf{Ejemplo de caso de uso}\\
		\toprule
		\textbf{Versión}              & 1.0    \\
		\textbf{Autor}                & Teodoro Ricardo García Sánchez \\
		\textbf{Requisitos asociados} & RF-xx, RF-xx \\
		\textbf{Descripción}          & La descripción del CU \\
		\textbf{Precondición}         & Precondiciones (podría haber más de una) \\
		\textbf{Acciones}             &
		\begin{enumerate}
			\def\labelenumi{\arabic{enumi}.}
			\tightlist
			\item Pasos del CU
			\item Pasos del CU (añadir tantos como sean necesarios)
		\end{enumerate}\\
		\textbf{Postcondición}        & Postcondiciones (podría haber más de una) \\
		\textbf{Excepciones}          & Excepciones \\
		\textbf{Importancia}          & Alta o Media o Baja... \\
		\bottomrule
	\end{tabularx}
	\caption{CU-1 Nombre del caso de uso.}
\end{table}




\apendice{Especificación de diseño}

\section{Introducción}

\section{Diseño de datos}
Para almacenar la información necesaria para la aplicación he diseñado una base de datos 
que sirva para facilitar el procesado de las reseñas y su posterior análisis 
por parte de los usuarios de la aplicación:

\imagen{Arquitectura Base de datos}{Arquitectura de la base de datos}

Como vemos almacenamos la información relativa a los conjuntos de datos en la tabla "datasets".
En esta tabla almacenamos la información necesaria para cargarlo como es el tipo (type) y el payload.
Según el tipo de conjunto de datos el método para cargarlo será diferente.
Una vez iniciado el proceso de carga, la información relativa a las reseñas se almacenan en la tabla "Review", "Products" y "ReviewUsers".

Por último, se almacena en la tabla "sentiments" la información relativa al análisis de sentimiento.
Se incluye la información obtenida del modelo así como el tiempo que ha tardado y el número de tokens procesados.
Este último campo (tokens) se puede utilizar para calcular el coste en el caso de OpenAI.
También se indica si la respuesta del modelo se considera errónea por tener un formato incorrecto. 
En este caso la respuesta del modelo se almacena en el campo "source".

\section{Diseño procedimental}

\imagen{Diagrama de secuencia login}{Login}
\imagen{Diagrama de secuencia admin}{Administración}
\imagen{Diagrama de secuencia gestor}{Gestor}


\section{Diseño arquitectónico}

\subsection{Arquitectura de la aplicación}

En el caso de la arquitectura de la aplicación, vemos que está centrada en la base de datos.
Desde el módulo de administración se pueden gestionar los usuarios y las reseñas.
Desde el módulo de servicios se cargan las reseñas y se gestionan los diferentes procesos de análisis.
A través de un servicio web se expone al usuario los diferentes informes de resultados.

\imagen{Arquitectura de la aplicación}{Arquitectura de la aplicación}

\subsection{Interfaz de usuario}
En la pantalla de login debe aparecer el logo de la aplicación 
junto con la información necesaria para poder acceder a la aplicación.
Se puede mostrar también una pequeña introdución sobre el proposito de la misma y 
un pequeño vínculo para acceder al registro como usuario.
\imagen{login}{Pantalla de login.}
En esta pantalla se mostrará la información relativa a las reseñas y 
al análisis de sentimiento realizado hasta el momento.
\imagen{Inicio}{Pantalla usuario gestor.}
En la pantalla de administración se podrá gestionar los usuarios y las reseñas.
\imagen{administración}{Pantalla de administración.}

\apendice{Documentación técnica de programación}

\section{Introducción}

En este apéndice se va a explicar cómo está organizado el proyecto, el
manual del programador y los requisitos necesarios para poder ejecutar el
proyecto.


\section{Estructura de directorios}

En el proyecto hay una estructura de directorios que intenta organizar el trabajo de manera lógica:
Por un lado tenemos el código y por otra la documentación.
Dentro del código podemos apreciar una carpeta llamada colab con un cuaderno de Jupyter para poder ejecutarlo en Colab.
Hay otra carpeta con el código del proyecto es sí.
La estructura de directorios quedaría así:

\begin{itemize}
	\item Código
	\begin{itemize}
		\item colab: Contiene un cuaderno de Jupyter preparado para ejecutar un ejemplo en Colab.
		\item prompt\_sentiment: Contiene el código del trabajo.
		\begin{itemize}
			\item frontend: Contiene el código de la interfaz de usuario. Está hecha con Vue, usando el plugin Vuetify.
			\item backend: Contiene el código del servicio web de la aplicación. Está hecha en Flask (python) usando Flask-restx y SQLAlchemy para el acceso a las base de datos.
			\item ngix: En este directorio está la configuración del servidor ngix que se usa en la aplicación como proxy inverso y para gestionar la configuración https.
		\end{itemize}
	\end{itemize}
	\item Documentación
	\begin{itemize}
		\item Manual de usuario: El manual de uso de la aplicación
		\item Memoria: Es la memoria del proyecto
	\end{itemize}
\end{itemize}

\section{Manual del programador}

\subsection {Instalación local}
Para la instalación local del proyecto tenemos realizar las siguientes instalaciones:

\subsubsection {Instalación de Python 3 ~\cite{lutz2001programming}}
Para ejecutar nuestro proyecto es necesario instalar Python, para ello
podemos descargar la versión aquí: \url{https://www.python.org/downloads/}
Cuando la descarga haya finalizado ejecutamos e instalamos.
También se recomienda usar conda para crear un entorno virtual para python. 
\url{https://docs.conda.io/projects/conda/en/latest/index.html}

\subsubsection {Instalación de node.js}
Para la parte de angular es necesario instalar Nodejs, por tanto accederemos 
a la página oficial, \url{https://nodejs.org/es/download/}, 
y descargaremos la versión que necesitemos y llevaremos a cabo su instalación

\subsubsection {Instalación de PostgreSql}
Para la instalación de PostgreSql tenemos que bajarnos la versión adecuada a nuestro sistema operativo e instalarla:
\url{https://www.postgresql.org/download/}

\subsection {Docker}
Para la instalación de docker sólo hay que instalar Docker-Desktop.
\subsubsection {Instalación de Docker Desktop}
\url{https://docs.docker.com/engine/install/}

\section{Compilación, instalación y ejecución del proyecto}
\subsection {Ejecución local}
El proyecto tiene 2 partes claramente diferenciadas: Frontend y Backend.
Para la ejecución del backend tenemos que abrir una consola de comandos, movernos al directorio backend 
y ejecutar:
//Instalamos las referencias
pip install -r requirements.txt

Creamos la base de datos (sólo en el caso de que no exista)
\begin{verbatim}
python manage.py crear
python manage.py rellenar (en caso de que esté vacía)
python manage.py run
\end{verbatim}

Para el frontend, nos cambiamos la carpeta Frontend y ejecutamos:
\begin{verbatim}
npm install
npm run serve
\end{verbatim}

\subsection {Ejecución en docker}
Para la ejecución en docker sólo hay que ejecutar en el directorio raiz del código:
\begin{verbatim}
docker-compose up
\end{verbatim}

\section{Pruebas del sistema}

Para la ejecución de las pruebas unitarias del backend hay que ejecutar:
\begin{verbatim}
pytest
\end{verbatim}
\apendice{Documentación de usuario}

\section{Introducción}

\section{Requisitos de usuarios}

\section{Instalación}

\section{Manual del usuario}



\apendice{Anexo de sostenibilización curricular}

\section{Introducción}
Este anexo incluirá una reflexión personal del alumnado sobre los aspectos de la sostenibilidad 
que se abordan en el trabajo.
Se pueden incluir tantas subsecciones como sean necesarias con la intención de explicar 
las competencias de sostenibilidad adquiridas durante el alumnado y aplicadas al Trabajo de Fin de Grado.

Más información en el documento de la CRUE 
\url{https://www.crue.org/wp-content/uploads/2020/02/Directrices_Sosteniblidad_Crue2012.pdf}.
Este anexo tendrá una extensión comprendida entre 600 y 800 palabras.

La sostenibilidad de un proyecto de software como el mencionado dependerá de varios factores, 
y es crucial considerar tanto los aspectos económicos como los sociales y ambientales. 
Aquí hay algunas reflexiones sobre la sostenibilidad del proyecto:

Económica:
Modelo de Negocio Sostenible: Un modelo de negocio sólido y sostenible es esencial. 
Evalúa la capacidad del proyecto para generar ingresos de manera constante 
y asegúrate de que el modelo de monetización elegido sea viable a largo plazo.
Gestión de Costos: Controlar eficientemente los costos de desarrollo, 
mantenimiento y operación es clave para la sostenibilidad económica. 
Considera la posibilidad de utilizar servicios en la nube de manera eficiente 
y evalúa constantemente la rentabilidad de las herramientas y tecnologías utilizadas.
Social:
Impacto Positivo: Evalúa cómo el proyecto puede tener un impacto positivo en la sociedad. 
Por ejemplo, si el software ayuda a las empresas a comprender y mejorar la satisfacción del cliente, 
esto podría contribuir a un impacto social positivo.
Accesibilidad: Asegúrate de que el software sea accesible para la mayor cantidad de personas posible. 
Considera aspectos como la usabilidad, la inclusión de personas con discapacidades 
y la disponibilidad en diferentes idiomas.
Ética en la Inteligencia Artificial: 
Si el proyecto utiliza inteligencia artificial (como OpenAI), 
es crucial considerar aspectos éticos. 
Asegúrate de que el uso de IA sea transparente, justo y ético, y evita sesgos o discriminación.
Ambiental:
Eficiencia Energética: Evalúa la eficiencia energética del software y las tecnologías utilizadas. 
La sostenibilidad ambiental implica minimizar el consumo de recursos, incluida la energía.
Impacto del Desarrollo: Considera el impacto ambiental durante el desarrollo del software. 
Esto podría incluir prácticas de desarrollo sostenibles, 
como la reducción de residuos de código o la eficiencia en el uso de recursos de desarrollo.
Comunidad y Colaboración:
Involucramiento de la Comunidad: Fomenta la participación y retroalimentación de la comunidad. 
Una comunidad activa puede contribuir al éxito continuo del proyecto 
y proporcionar ideas valiosas para mejoras sostenibles.
Colaboración Abierta: Considera la posibilidad de colaborar con otros 
proyectos de código abierto y contribuir al ecosistema más amplio. 
La colaboración puede generar sinergias y mejorar la sostenibilidad a largo plazo.
En resumen, para lograr la sostenibilidad de un proyecto de software, 
es necesario equilibrar aspectos económicos, sociales y ambientales. 
La transparencia, la ética y la consideración de los impactos a largo plazo 
son esenciales para construir y mantener un proyecto que beneficie tanto 
a los usuarios como al entorno en el que se desarrolla.


\bibliographystyle{plain}
\bibliography{bibliografiaAnexos}

\end{document}
